\begin{textsample}{POS Dim 2 – pe – Score 22.00 – pe/pe\_em\_29.txt}  \label{ex:f2_pos_273}
3.5 – O \textbf{ITINERÁRIO} FORMATIVO DE EDUCAÇÃO PROFISSIONAL \textbf{TÉCNICA} Nos termos da Lei no 13.415/2017, as áreas do conhecimento, assim como a Formação \textbf{Técnica} e Profissional, constituem a flexibilização da etapa do Ensino Médio.
A educação profissional e tecnológica não está ligada apenas às tecnologias educacionais e à educação \textbf{técnica} de forma isolada.
Ela se coaduna à formação interdimensional do sujeito em competências tanto cognitivas quanto socioemocionais.
Isso porque, diante de intensas mudanças globais e do surgimento constante de novas tecnologias, a escola e o educando são desafiados continuamente a repensarem comportamentos e atitudes no cumprimento da missão de impactar o espaço social com projetos e ações inovadoras e sustentáveis.
Por essa razão, diante das novas formas de produção e das dinâmicas sociais, a educação profissional e tecnológica deve fomentar o olhar crítico e reflexivo dos sujeitos, buscando, dessa forma, o compromisso com uma formação integrada e abrangente para o mundo do trabalho e da tecnologia.
Além disso, as demandas sociais e educacionais contemporâneas são imperativas quanto à necessidade de estruturar um \textbf{itinerário} formativo \textbf{flexível}, dinâmico e abrangente, segundo os interesses dos indivíduos e possibilidades das \textbf{instituições} educacionais.
Como já mencionado, os \textbf{itinerários} formativos das diferentes áreas e da formação \textbf{técnica} e profissional devem ser organizados considerando os quatro eixos estruturantes a seguir : investigação científica, processos criativos, mediação e intervenção sociocultural e empreendedorismo.
Ainda no \textbf{tocante} ao \textbf{itinerário} de formação \textbf{técnica} e profissional, o parágrafo 14o do \textbf{Artigo} 12 das DCNEM traz, “ [... ] deve observar a integralidade de ocupações \textbf{técnicas} reconhecidas pelo setor produtivo, tendo como referência a Classificação Brasileira de Ocupações ( CBO ) ”. ( BRASIL, 2018 ).
Sendo assim, no \textbf{itinerário} formativo de educação profissional e \textbf{técnica} é proposta, em grande medida, uma articulação curricular entre o ensino propedêutico e a formação \textbf{técnica} e profissional.
Nesse processo, fazem -se presentes também o fomento ao protagonismo juvenil, projeto de vida e, consequentemente, a \textbf{garantia} do direito de escolha do estudante, bem como um percurso formativo potencialmente voltado ao desenvolvimento de competências e habilidades conectadas às demandas de formação do profissional do futuro.
O \textbf{itinerário} de formação \textbf{técnica} e profissional, segundo o parágrafo único, do inciso X, do \textbf{artigo} 6.o, do das DCNEM, compreende um conjunto de termos e conceitos próprios, tais como : a ) ambientes simulados : são ambientes pedagógicos que possibilitam o desenvolvimento de atividades práticas da aprendizagem profissional quando não puderem ser elididos riscos que sujeitem os aprendizes à insalubridade ou à periculosidade nos ambientes reais de trabalho ; b ) formações experimentais : são formações autorizadas pelos respectivos sistemas de ensino, nos termos de sua \textbf{regulamentação} específica, que ainda não constam no \textbf{Catálogo} Nacional de \textbf{Cursos} \textbf{Técnicos} ( CNCT ) ; c ) aprendizagem profissional : é a formação técnico-profissional compatível com o desenvolvimento físico, moral, psicológico e social do jovem, de 14 a 24 anos de idade, \textbf{previsto} no § 4o do \textbf{art}. 428 da Consolidação das Leis do Trabalho ( CLT ) e em \textbf{legislação} específica, caracterizada por atividades teóricas e práticas, metodicamente organizadas em tarefas de complexidade progressiva, conforme respectivo perfil profissional ; d ) \textbf{qualificação} profissional : é o processo ou resultado de formação e desenvolvimento de competências de um determinado perfil profissional, definido no mercado de trabalho ; e ) \textbf{habilitação} profissional \textbf{técnica} de nível médio : é a \textbf{qualificação} profissional formalmente reconhecida por meio de diploma de conclusão de \textbf{curso} \textbf{técnico}, o qual, quando registrado, tem validade nacional ; f ) programa de aprendizagem : compreende arranjos e combinações de \textbf{cursos} que, articulados e com os devidos aproveitamentos curriculares, possibilitam um \textbf{itinerário} formativo.
A \textbf{oferta} de programas de aprendizagem tem por objetivo apoiar trajetórias formativas, que tenham relevância para os jovens e favoreçam sua inserção futura no mercado de trabalho.
Observadas as normas \textbf{vigentes} relacionadas à \textbf{carga} \textbf{horária} mínima e ao tempo máximo de duração do contrato de aprendizagem, os programas de aprendizagem podem compreender distintos arranjos ; g ) \textbf{certificação} intermediária : é a possibilidade de emitir \textbf{certificação} de \textbf{qualificação} para o trabalho quando a formação for estruturada e organizada em etapas com terminalidade ; h ) \textbf{certificação} profissional : é o processo de avaliação, reconhecimento e \textbf{certificação} de saberes adquiridos na educação profissional, inclusive no trabalho, para fins de \textbf{prosseguimento} ou conclusão de estudos nos termos do \textbf{art}. 41 da LDB ( Brasil, 2018 ).
As trilhas de EPT sugeridas em Pernambuco- em anexo- atendem, na sua concepção, aos princípios norteadores constantes no organizador curricular do \textbf{itinerário} formativo de educação profissional e \textbf{técnica}.

% matched lemmas: art, artigo, carga, catálogo, certificação, curso, flexível, garantia, habilitação, horário, instituição, itinerário, legislação, oferta, prever, prosseguimento, qualificação, regulamentação, tocante, técnico, vigente
\end{textsample}
